\chapter{Requirements} 
\addcontentsline{toc}{chapter}{Requirements} % rucne pridanie do obsahu
\markboth{Introduction}{Requirements} % vyriesenie hlaviciek

\section{Application}

\subsection{Client -- Server}

The system architecture is based on the client--server model.
It is recommended to avoid hardcoded constants and instead allow the user to configure the relevant parameters through a configuration file or interface.

\subsection{Client}

The client component should be as simple and lightweight as possible.
Its main functions are:

\begin{itemize}
    \item detecting the operating system version, ideally including installed patches,
    \item compiling a list of installed applications and their versions,
    \item sending the collected data to the server, including the IP address, MAC address, and computer name,
    \item cross-platform support -- Windows (XP through 11), Linux (Ubuntu), macOS, Android, and iOS.
\end{itemize}

A separate consideration is the distribution of the client to endpoint devices and the configuration of communication with the server.

\subsection{Server}

The server component provides the following functions.

\subsubsection{Data Collection from Clients}

The server receives and stores data sent by individual clients.

\subsubsection{IP Address Scanning}

The server scans the network to detect devices that are connected but do not have the client installed.
Possible approaches include:

\begin{itemize}
    \item scanning a defined IP address range with the option to exclude addresses where clients are running,
    \item listening to DHCP or DNS protocols and comparing the source IP address against the list of registered clients,
    \item using Suricata or a similar tool to perform network scanning.
\end{itemize}

\subsubsection{Client Display}

The following information is displayed for each client:

\begin{itemize}
    \item the logical name of the device and its IP address,
    \item the time of last connection (\textit{last seen}); a trigger is fired when a configured limit is exceeded (e.g.\ no communication for 7 days); a lightweight alarm is also activated if the client is offline but the device is detected on the network (which may indicate that the client was uninstalled),
    \item the client's risk score displayed as a change indicator (e.g.\ arrow up, dash, arrow down).
\end{itemize}

\paragraph{Device Grouping}

Devices can be assigned to groups with an associated risk coefficient.
The following group types are distinguished:

\begin{itemize}
    \item \textbf{Hardware groups} -- servers, desktop computers, laptops, mobile devices, and others.
    \item \textbf{Organizational groups} (without a risk coefficient) -- HR, management, finance department, developers, or location.
    \item \textbf{Groups by location and priority} (with a risk coefficient) -- perimeter, critical zone, standard internal department.
\end{itemize}

Assigning devices to groups can be done manually or automatically (e.g.\ based on network range).
Reassigning devices between groups is done manually.
The system also allows removing a device from the list, for example when it is physically decommissioned.

Each group displays a risk score with a change indicator (arrow up, dash, arrow down).

\paragraph{Finding Records}

Each finding is assigned a unique identifier (e.g.\ \texttt{ID345-as23}), enabling integration with ticketing systems.
The system maintains a historical view of patch development within the organization, for example: \texttt{CVE-XXX-XXX} -- device, date from--to, finding ID.
The log retention period is configurable; the default value is set to \textbf{13 months}.

\paragraph{Allowed Application List}

The system allows maintaining a list of allowed applications. An administrator can add an application to the whitelist by clicking on it directly in the list of installed applications on a given device.

\subsubsection{Dashboards and Statistics}

The server provides the following overviews:

\begin{itemize}
    \item \textbf{Device statistics} -- total number of devices including those without the client installed, counts per group.
    \item \textbf{Risk overviews} -- overall organizational risk and its development over time (graph), graphs for individual groups.
    \item \textbf{Prohibited applications} -- occurrences of prohibited applications and their frequency.
\end{itemize}

\subsubsection{Vulnerability Database Communication}

The server communicates with external sources to retrieve:

\begin{itemize}
    \item CVSS scores for identified vulnerabilities,
    \item additional information, ideally including an indicator of whether the vulnerability is being actively exploited.
\end{itemize}

% ============================================================
\section{Risk Assessment Methodology}

This section describes a rough outline of the risk assessment methodology with a minimal number of predefined constants.

\subsection{Input Parameters}

\begin{itemize}
    \item \textbf{CVSS score} -- the numerical severity of a vulnerability, e.g.\ $\text{CVSS} = 9{.}5$.
    \item \textbf{Application criticality for the organization} -- coefficient $k_a$ with a default value of $k_a = 1$:
    \begin{itemize}
        \item critical applications (OS, EDR, and similar; may be predefined with the option to extend): $k_a = 2$,
        \item other applications (manual configuration possible): $k_a = 1$.
    \end{itemize}
    \item \textbf{Whitelist coefficient} -- coefficient $k_w$ expressing the application's whitelist status:
    \begin{itemize}
        \item allowed application: $k_w = 1$,
        \item prohibited application: $k_w = 3$.
    \end{itemize}
\end{itemize}

\subsection{Application Risk Calculation for a Device}

The application risk for a given device is calculated as:

\begin{equation}
    R_{\text{app}} = \sum_{i=1}^{n} \left( \text{CVSS}_i \cdot k_{a,i} \cdot k_{w,i} \right) + N_{\text{banned}} \cdot 3,
    \label{eq:risk_app}
\end{equation}

where:
\begin{itemize}
    \item $n$ is the number of applications with identified vulnerabilities,
    \item $\text{CVSS}_i$ is the score of the $i$-th vulnerability,
    \item $k_{a,i}$ is the criticality coefficient of the $i$-th application,
    \item $k_{w,i}$ is the whitelist coefficient of the $i$-th application,
    \item $N_{\text{banned}}$ is the number of prohibited applications with no recorded vulnerability.
\end{itemize}

\subsection{Groups and Criticality Coefficients}

Each device belongs to a group with an assigned criticality coefficient $k_g$:

\begin{table}[h!]
    \centering
    \begin{tabular}{lc}
        \hline
        \textbf{Group Type} & \textbf{Coefficient $k_g$} \\
        \hline
        Extra critical group & $5$--$7$ \\
        Perimeter            & $5$ \\
        Critical group       & $3$ \\
        Standard group       & $1$ \\
        \hline
    \end{tabular}
    \caption{Group criticality coefficients.}
    \label{tab:group_coeff}
\end{table}

\subsection{Device and Group Risk Calculation}

The risk of an individual device is calculated as:

\begin{equation}
    R_{\text{PC}} = R_{\text{app}} \cdot k_g,
    \label{eq:risk_pc}
\end{equation}

where $k_g$ is the criticality coefficient of the group to which the device belongs (see Table~\ref{tab:group_coeff}).

The total risk of a group is the sum of the risks of all devices within that group:

\begin{equation}
    R_{\text{group}} = \sum_{j=1}^{m} R_{\text{PC},j},
    \label{eq:risk_group}
\end{equation}

where $m$ is the number of devices in the group.

\subsection{Risk Level Thresholds}

Thresholds for risk levels (green, yellow, red) must be defined to evaluate the overall state of the organization.
These thresholds are subject to further calibration.

An important rule is that \textbf{even if the overall risk value falls within the green zone, the presence of a critical issue on the perimeter automatically prevents a green overall rating}.
Formally, this rule can be expressed as:

\begin{equation}
    \text{OverallRisk} =
    \begin{cases}
        \text{red,} & \text{if } R_{\text{perimeter}} > \theta_{\text{perimeter}}, \\
        f\!\left(\sum_{\text{groups}} R_{\text{group}}\right), & \text{otherwise,}
    \end{cases}
    \label{eq:overall_risk}
\end{equation}

where $\theta_{\text{perimeter}}$ is a configurable threshold for the perimeter and $f(\cdot)$ is a function mapping the numerical score to a color-coded risk level.